% GENERATED FILE. Edit canonical lessons, then run npm run generate:books.
% canonical-source-hash: fnv1a64:92aecb01d89fccc8
% canonical-lessons: GE-C06-sechs, GE-C06-sieben, GE-C06-acht

\chapter{Six, Seven, Eight}
\label{ch:ge-zahlen-6-8}
\hlchaptermodality{8}

\begin{chapteropening}
\textbf{By the end of this chapter:} \emph{I can count six, seven and eight in German, and hear the three sounds that make them hard to guess from the spelling.}

Three more numbers, each with the sound that makes it unguessable from the page: the chs of sechs, the buzzing s of sieben, and the throat sound in acht that English spells gh.
\end{chapteropening}

\section[sechs]{sechs — \textquotedblleft{}six,\textquotedblright{} where \emph{chs} is just an \emph{x}}

\label{lesson:GE-C06-sechs}



\begin{quote}
The second hand starts here. And the sixth number quietly settles the last thing German \textbf{ch} can do.
\end{quote}

\subsection*{You'll want to know: sechs}
\begin{quote}
\textbf{sechs} — \textquotedblleft{}six.\textquotedblright{}
\end{quote}

\begin{sounds}
\begin{itemize}
\raggedright
  \item \texttt{ch-none} — in \textbf{chs}, the \emph{ch} makes no throat sound at all. The pair is a plain \textbf{ks}, exactly the English letter \emph{x}.
\end{itemize}

So \emph{sechs} is \emph{zeks}, and it rhymes with English \emph{sex} — which is also, as it happens, the Latin word for six.

That completes the \textbf{ch}: raspy after \emph{a} as in \emph{machen}, and silent-into-\emph{ks} when an \emph{s} follows. Nothing else to learn about it.
\end{sounds}

\begin{cousinweb}
\textbf{sechs} is English \textbf{six}, and here the two languages have barely moved apart at all — \emph{zeks} against \emph{siks}. Even the Latin cousin \emph{sex} is close enough to hear, because this is one of the few numbers no sound law rearranged much.

Which makes it the cheapest number in the chapter: if you can say \emph{x}, you can say \emph{sechs}.
\end{cousinweb}

\subsection*{Your turn}
Say these aloud:

\begin{itemize}
\raggedright
  \item \textquotedblleft{}sechs\textquotedblright{} — \emph{zeks}, no throat rasp
  \item \textquotedblleft{}machen\textquotedblright{} then \textquotedblleft{}sechs\textquotedblright{} — the two jobs of \emph{ch}
  \item \textquotedblleft{}sechs\textquotedblright{} then \textquotedblleft{}six\textquotedblright{} — barely a step apart
\end{itemize}

\begin{tcolorbox}[breakable,colback=teal!4,colframe=teal!35!black,title={Before you move on}]
What is \textquotedblleft{}six\textquotedblright{}? (\emph{Sechs}, said \emph{zeks}.) What does \textbf{ch} do in front of an \emph{s}? (Nothing — the pair is a plain \emph{ks}.) And after an \emph{a}? (The raspy throat sound, as in \emph{machen}.)
\end{tcolorbox}
\section[sieben]{sieben — \textquotedblleft{}seven,\textquotedblright{} and the buzzing \emph{s}}

\label{lesson:GE-C06-sieben}



\begin{quote}
You have been saying \emph{sechs} as \emph{zeks} without being told why. Here is the rule, on the number that needs it most.
\end{quote}

\subsection*{You'll want to know: sieben}
\begin{quote}
\textbf{sieben} — \textquotedblleft{}seven.\textquotedblright{} Two syllables: \emph{ZEE-ben}.
\end{quote}

\begin{sounds}
\begin{itemize}
\raggedright
  \item \texttt{s-voiced} — a German \textbf{s} in front of a vowel \textbf{buzzes}, like the English letter \emph{z}. At the end of a word it goes back to a plain hiss.
  \item \texttt{ie-long-ee} — the long \emph{ee} again, as in \emph{vier}.
\end{itemize}

So \emph{sieben} is \emph{ZEE-ben}, not \emph{SEE-ben}. And it explains the two numbers before it: \emph{sechs} is \emph{zeks} and not \emph{seks} for the same reason.
\end{sounds}

\begin{cousinweb}
\textbf{sieben} is English \textbf{seven}, with the middle consonant softened from a \emph{v} to a \emph{b} — a swap those two letters make in each other's direction constantly, because they are made with the same lips.

This is the one number where the Latin cousin is also unmistakable: \emph{septem}, which English still uses in \emph{September} and in \emph{septet}. Seven is the number almost nobody's language disguised.
\end{cousinweb}

\subsection*{Your turn}
Say these aloud:

\begin{itemize}
\raggedright
  \item \textquotedblleft{}sieben\textquotedblright{} — \emph{ZEE-ben}, buzz the \emph{s}
  \item \textquotedblleft{}sechs, sieben\textquotedblright{} — both start with the buzz
  \item \textquotedblleft{}sieben\textquotedblright{} then \textquotedblleft{}seven\textquotedblright{}
\end{itemize}

\begin{tcolorbox}[breakable,colback=teal!4,colframe=teal!35!black,title={Before you move on}]
What is \textquotedblleft{}seven\textquotedblright{}? (\emph{Sieben}, said \emph{ZEE-ben}.) What does a German \textbf{s} do in front of a vowel? (It buzzes, like \emph{z}.) And at the end of a word? (It hisses.) Which English word is \emph{sieben}? (\emph{Seven} — the \emph{v} softened to a \emph{b}.)
\end{tcolorbox}
\section[acht]{acht — \textquotedblleft{}eight,\textquotedblright{} on a rasp you already have}

\label{lesson:GE-C06-acht}



\begin{quote}
Nothing new in the mouth for this one. You met its sound on the very first \emph{Gute Nacht}.
\end{quote}

\subsection*{You'll want to know: acht}
\begin{quote}
\textbf{acht} — \textquotedblleft{}eight.\textquotedblright{} Said \emph{akht}.
\end{quote}

The \textbf{ch} is the raspy back-of-the-throat sound from \emph{Nacht} and \emph{machen}: a clean open \emph{ah}, then the scrape. One syllable, no vowel after it.

\begin{cousinweb}
\textbf{acht} is English \textbf{eight} — and the pairing is easier to believe once you notice that English still \emph{writes} the missing sound.

\emph{Eight} has a \textbf{gh} in it that nobody says. It is not decoration. English used to make a real throat scrape there, exactly the one German still makes, and the spelling is the fossil left behind after the sound fell out.

So German \textbf{-cht-} answers English \textbf{-ght-}, and it does it every single time:

\noindent
\begin{tabularx}{\linewidth}{@{}>{\raggedright\arraybackslash}X>{\raggedright\arraybackslash}X@{}}
\toprule
\textbf{German} & \textbf{English} \\
\midrule
a\textbf{cht} & ei\textbf{ght} \\
Na\textbf{cht} & ni\textbf{ght} \\
Li\textbf{cht} & li\textbf{ght} \\
re\textbf{cht} & ri\textbf{ght} \\
\bottomrule
\end{tabularx}

\emph{Nacht} and \emph{night} you already own. The pattern is worth more than the number: whenever you see a German word with \textbf{cht}, try the English one with \textbf{ght}.
\end{cousinweb}

\subsection*{Your turn}
Say these aloud:

\begin{itemize}
\raggedright
  \item \textquotedblleft{}acht\textquotedblright{} — \emph{akht}, open \emph{ah} then the scrape
  \item \textquotedblleft{}Nacht\textquotedblright{} then \textquotedblleft{}acht\textquotedblright{} — the same ending twice
  \item the pairs — \textquotedblleft{}acht/eight\textquotedblright{}, \textquotedblleft{}Nacht/night\textquotedblright{}, \textquotedblleft{}Licht/light\textquotedblright{}
\end{itemize}

\begin{tcolorbox}[breakable,colback=teal!4,colframe=teal!35!black,title={Before you move on}]
What is \textquotedblleft{}eight\textquotedblright{}? (\emph{Acht}.) What does German \textbf{-cht-} answer in English? (\textbf{-ght-}.) Why does \emph{eight} have a silent \textbf{gh}? (English used to say the same scrape German still says.) Give another pair. (\emph{Nacht}/\emph{night}.)
\end{tcolorbox}
